\section{Threats to Validity}\label{sec:threats}

\textbf{Model access.} The models were reached through different routes (Table~\ref{tab:models}). This is a known risk for studies that use LLMs \citep{sallou2024breaking}. Three were used through conversational interfaces rather than a raw API. Claude ran through Claude Code subagents, whose harness adds its own system prompt to ours. GPT and DeepSeek ran through their web chats, whose own system prompts and defaults are outside the study's control. So temperature and the output limit were controlled only for Gemini and Qwen. This affects the consistency comparison of RQ3 more than any other result. GPT was generated with ChatGPT's default model, so its exact version is not known. It answered with a medium reasoning effort, left unchanged. One author (T.~Hussain) operated its web sessions following a written protocol: a new temporary, unpersonalised chat per prompt, no tools, no follow-up, and replies copied unedited. Adherence to this protocol was not independently observed. For 16 of its 42 prompts, the interface turned the pasted prompt into an attached file with identical content. They came from the three projects with the longest documentation. Gemini is a flash-tier model rather than the planned Pro tier. Its hidden reasoning tokens count against its 8,192-token output limit. Qwen's free tier rejects the prompts of the two largest projects, so its missing cells reflect access, not quality. A sixth model, Llama, was dropped before the main run because no working access route was found.

\textbf{Training-data contamination.} All projects are public, and only two of the seven manifest repositories have fewer than 2,000 GitHub stars. So a model may reproduce a memorised manifest rather than reason from the documentation. We record provider refusals for recitation. We also compare each model's similarity to the human manifest with its similarity to the other models' configurations for the same project. Only an excess over this peer level is read as reproduction (Section~\ref{sec:disc-memorisation}). Reproduction would pull a model's configurations toward the human manifest, and so shrink the differences reported for RQ1.

\textbf{Selection of projects and of the human baseline.} Results describe these seven projects, not application classes in general (Section~\ref{sec:disc-generalisation}). The set was fixed after a two-project pilot and before the main run. A project was selected only if its human manifest deployed on the evaluation cluster. So the human D5 outcomes are partly true by construction, and the selection favours applications that run on a local cluster. Human runtime results are therefore reported as a reference only.

\textbf{Evaluation cluster.} D5 and D6 are measured on a four-node \texttt{kind} cluster on a single host. Cloud load balancers, multi-zone scheduling and production storage are not exercised. Absolute times should not be read as representative of production. \texttt{kind}'s local-path storage does not apply a pod's \texttt{fsGroup}, which could make a database fail here that would start elsewhere. None of the eight database failures depends on it. Seven run the container as root with every capability dropped, which prevents the entrypoint from changing file ownership or switching user, whatever the volume's group. The eighth fails because another user owns its data directory, and \texttt{fsGroup} does not change that.

\textbf{Measurement adjustments.} The pipeline was corrected 25 times during the study (Appendix~\ref{app:fixes}). Each correction was logged with its trigger, applied to every configuration it could affect, and followed by re-evaluation. All were made after generation had begun. Most were found while diagnosing failures, so they are one-sided. Of the corrections, 8 affected only the running of the pipeline and 3 made scoring stricter. The other 14 removed a failure that the investigation traced to the measurement or to an undocumented prerequisite. They never removed a success that turned out to be spurious. Appendix~\ref{app:fixes} separates corrections of the measurement from changes to an outcome's definition and from supplied prerequisites. Among the definition changes that removed failures, Correction~19 leaves a strict variant of D6, which is reported alongside. Correction~26 changes the outcome of two configurations (DeepSeek, Grafana, production prompt). Without it, the production smoke-test rate would be 58 rather than 60 of 99. Two corrections (2 and 10) were triggered by failures of human manifests. The deployment results are robust to leaving out any one project (Section~\ref{sec:res-rq5}). Two adjustments supply facts that the bundles lack (image aliases and start-up prerequisites, Section~\ref{sec:method-pipeline}). They apply to 29 generated configurations and to no human manifest. Take the worst case, in which every adjusted configuration that succeeded would otherwise have failed. Then the share of generated configurations both fully Ready and passing the smoke test falls from 59.6\% to 51.0\%. Four of the dashboard failures concern a start-up Secret that the documentation never mentions. Once the Secret is supplied, these configurations fail because their RBAC rules do not let the dashboard update it. A deployer would also have to discover this. Three configurations set a required variable to an empty string, which the application rejects. The start-up adjustment fires only when a variable is absent, so they are scored as failing.

\textbf{Labelling and tools.} Root-cause labels were assigned by a single rater. They were drafted with the help of an AI assistant from the same model family as one of the evaluated models (Claude). The blind second labelling pass was also made by a Claude model, not by a second human. It shows that the labels follow from the recorded evidence ($\kappa=0.84$). But it cannot rule out errors that models of one family share. The evidence level of each label is reported. Every label is published with its evidence so that it can be re-rated. The static tools are imperfect and disagree with one another. So we read their findings as what these pinned tool versions detect. We report the agreement between the two security scanners. Redaction decisions followed a written rule set. Every redaction was logged and reviewed by a second author, but borderline cases remain possible. Facts that the bundles do not contain are excluded from D4. However, the default ports of widely used database images are treated as general knowledge and scored.

\textbf{Resilience measurement.} Kubernetes controllers handle the two injected faults whatever the configuration says. So recovery success barely discriminates, and downtime depends mainly on replica count. D6 therefore characterises the replica choices a model makes, not a broader notion of resilience. Faults such as node drains or resource pressure would test more of a configuration's choices.

\textbf{Drift.} Providers update the weights behind a model identifier without notice. An unpinned image tag resolves to whatever is current. Generation dates are recorded for every configuration, and images were resolved in November 2026. Tag drift affects the human baselines too. The human Sock Shop manifest references \texttt{mongo} without a tag. It now resolves to a release whose wire protocol the project's 2017 cart service no longer speaks. Re-deploying with the image pinned to an older release passed, but the baseline is scored as published.
